problem:  
Let \((M,\Delta,g_\Delta)\) be a complete, equiregular, sub‑Riemannian manifold of bounded geometry, with Popp measure \(\mu\).  
For a fixed \(1\le k<\dim(\Delta)\), define the large‑scale filling exponent
\[
\beta(k,M) = \limsup_{r\to\infty} \frac{\log \mathrm{Fill}_M^{(k)}(r)}{\log r},
\]
where \(\mathrm{Fill}_M^{(k)}(r)\) is the \(k\)-dimensional Gromov filling function defined via horizontal integral currents (Ambrosio–Kirchheim framework).

At each point \(p\in M\), let the nilpotent approximation (tangent Carnot group) be \(\mathrm{Nil}_p(M)\) with homogeneous dimension \(Q(p)\), and suppose  
\[
\beta(k,p):=\limsup_{r\to\infty}\frac{\log \mathrm{Fill}_{\mathrm{Nil}_p(M)}^{(k)}(r)}{\log r}
\]
exists and is finite.  

**Problem:**  
Assuming that the map \(p\mapsto \mathrm{Nil}_p(M)\) varies continuously in the Gromov–Hausdorff sense and that the set \(\{\beta(k,p)\}\) attains finitely many values, determine whether
\[
\beta(k,M) = \sup_{p\in M} \beta(k,p)
\]
must hold.  
If not, classify possible mechanisms leading to strict inequality in terms of:
* nonuniformity or curvature of the horizontal distribution,
* topological obstructions (e.g., nontrivial holonomy or abnormal geodesics),
* or large‑scale geometric anisotropies not detected by local tangent cones.

Equivalently, is the global asymptotic filling complexity of \(M\) completely governed by its **worst local Carnot model**, or can global geometry increase (or decrease) filling difficulty beyond any local tangent model?

Why is it a "good" problem:  
This refined formulation poses a sharp, invariant comparison between global and local filling behavior without relying on ad hoc constants or averaging.

- **Profound insight:**  
It seeks a foundational structural link between infinitesimal Carnot geometries (local nilpotent models) and global geometric‑measure invariants (filling exponents). Establishing or refuting equality would uncover how sub‑Riemannian curvature and topology alter large‑scale filling phenomena beyond tangent‑cone data.
  
- **Bridge between fields:**  
The problem integrates **sub‑Riemannian geometry** (nilpotentization, Popp measure) with **geometric measure theory** (currents and fillings) and **asymptotic geometry** (behavior under scaling and Gromov–Hausdorff convergence). Any progress would require blending analytic, geometric, and group‑theoretic tools.
  
- **Extreme simplicity:**  
In essence, the question asks:
> “Does the hardest tangent Carnot geometry already dictate the global filling exponent of the manifold?”
This concise statement hides deep subtleties of non‑Riemannian curvature and global topology, positioning it as a natural and potentially field‑defining conjecture in sub‑Riemannian asymptotic geometry.